% !TeX root = ../sections/02_exercises.tex
\begin{exercise}
	$n\geq 2$ のとき，不等式
	\[
	\frac{2^n}{n!}\leq 2\left(\frac{2}{3}\right)^{n-2}
	\]
	が成立することを証明し，極限
	\[
	\lim_{n\to\infty}\frac{2^n}{n!}
	\]
	を求めよ。
	ただし，$\epsilon N$ 論法を用いる必要はない。
\end{exercise}
\begin{background}
	この問題では，階乗 $n!$ が指数 $2^n$ よりも速く大きくなることを，
	具体的な不等式によって確認する。
	
	\begin{remark}
		左辺は
		\[
		\frac{2^n}{n!}
		=
		\frac{2}{1}\cdot\frac{2}{2}\cdot\frac{2}{3}
		\cdots\frac{2}{n}
		\]
		と分解できる。
		$n\geq 2$ のとき，$3$ 以降の因子について
		\[
		\frac{2}{k}\leq \frac{2}{3}
		\qquad (k\geq 3)
		\]
		が成り立つ。
	\end{remark}
	
	\subsection{階乗を積の形で評価する準備}
	
	\begin{lemma}[はさみうちの原理]
		数列 $\{x_n\},\{y_n\}$ が
		\[
		0\leq x_n\leq y_n
		\]
		を満たし，
		\[
		\lim_{n\to\infty}y_n=0
		\]
		であるならば，
		\[
		\lim_{n\to\infty}x_n=0
		\]
		である。
	\end{lemma}
	
	\begin{remark}
		$0<2/3<1$ なので，
		\[
		\left(\frac{2}{3}\right)^{n-2}\to 0
		\]
		である。
	\end{remark}
\end{background}
\begin{strategy}
	まず
	\[
	\frac{2^n}{n!}
	\]
	を積の形に直す。
	
	実際，
	\[
	\frac{2^n}{n!}
	=
	2\cdot 1\cdot
	\frac{2}{3}\cdot\frac{2}{4}
	\cdots
	\frac{2}{n}
	\]
	である。
	
	ここで $k\geq 3$ ならば
	\[
	\frac{2}{k}\leq\frac{2}{3}
	\]
	であるから，
	\[
	\frac{2^n}{n!}
	\leq
	2\left(\frac{2}{3}\right)^{n-2}
	\]
	が得られる。
	
	さらに
	\[
	0<\frac{2}{3}<1
	\]
	であるから，
	\[
	2\left(\frac{2}{3}\right)^{n-2}\to 0
	\]
	である。
	
	また
	\[
	0\leq \frac{2^n}{n!}
	\]
	なので，はさみうちの原理より
	\[
	\lim_{n\to\infty}\frac{2^n}{n!}=0
	\]
	である。
\end{strategy}